\section{Introduction}
\label{sec:introduction}

For positive integers $r$ and $k$, let $W_r(k)$ be the least positive
integer $N$ such that every map $[N]\to[r]$ is constant on some
arithmetic progression
\[
 a,a+d,\ldots,a+(k-1)d,\qquad a,d\ge1,\quad a+(k-1)d\le N.
\]
Here $[N]=\{1,\ldots,N\}$, and using fewer than $r$ colors is allowed.
The finiteness of these numbers is van der Waerden's theorem
\cite{vanDerWaerden1927}.
The lower-bound problem asks how long an interval can be colored while
avoiding a monochromatic progression of a prescribed length. We prove
a bound with one threshold on the length, valid for every number of
colors.

\begin{theorem}
\label{main:intro}
There is an absolute integer $K_0$ such that, for every integer
$k\ge K_0$ and every integer $r\ge2$,
\[
 W_r(k)>k^{c k\lfloor\log_2r\rfloor},
 \qquad c=10^{-5}.
\]
\end{theorem}

In particular, for each fixed $r\ge2$,
\[
 \lim_{k\to\infty}W_r(k)^{1/k}=\infty.
\]
Erd\H{o}s explicitly asked for this full limit in the two-color case
\cite[p.~90, (2)]{erdos1980}.
Thus Theorem~\ref{main:intro} gives a quantitative positive resolution
of his superexponential-growth question.
It is stronger than the assertion that $W_2(k)/2^k$ tends to infinity.

\paragraph{Earlier bounds.}
Early lower bounds were obtained by Erd\H{o}s and Rado
\cite{erdosRado1952} and by Schmidt~\cite{schmidt1962}.
Berlekamp's algebraic construction gives
$W_2(p+1)>p2^p$ for every prime $p$ \cite[Theorem~2]{berlekamp1968}.
The local lemma of Erd\H{o}s and Lov\'asz~\cite{erdosLovasz1975}
provided another route to strong coloring bounds.
Szab\'o~\cite[Theorem~2]{szabo1990} proved
$W_2(k)\ge 2^k/k^\varepsilon$ for every fixed $\varepsilon>0$
and all sufficiently large $k$.
Kozik and Shabanov~\cite[Theorem~2]{kozikShabanov2016}
subsequently obtained an absolute $\beta>0$ such that
$W_r(k)\ge\beta r^{k-1}$ for all $r\ge2$ and $k\ge3$.
Hunter~\cite[Theorem~1]{hunter2025} improved the exponential base
for each fixed $r\ge5$, with a threshold allowed to depend on $r$.

Recent results distinguish the two questions posed by Erd\H{o}s.
Fox and Hunter \cite[Theorem~1]{foxHunter2026} proved
$W_3(k)>2^{k\log^*k/4}$ for all sufficiently large $k$, where
$\log^*k$ counts the iterations of the natural logarithm needed to
reach a value at most one.
Campos, Fox, and Schildkraut \cite[Theorem~1]{camposFoxSchildkraut2026}
proved
\[
 W_2(k)\ge(1-o(1))k2^{k-1}.
\]
They credit ChatGPT 5.6 Sol Pro with the coloring and an initial proof.
Their result settles the question about $W_2(k)/2^k$, whereas
Theorem~\ref{main:intro} gives the full superexponential limit.
For every fixed $\varepsilon\in(0,1)$, Fox and Hunter also prove
$W_r(k)\ge r^{(1-\varepsilon)k\log k}$ when $k$ is sufficiently
large in terms of $\varepsilon$ and
$r\ge(\log k)^{3/\varepsilon}$ \cite[Theorem~2]{foxHunter2026}.
For example, taking $\varepsilon=1/2$ gives a stronger lower bound
when $r\ge(\log k)^6$ and $k$ is sufficiently large. The bound here
holds uniformly down to two colors.

Growth in the number of colors, with the progression length fixed,
has a different history. Behrend's construction and Rankin's extension
to longer progressions supply dense progression-free sets
\cite{behrend1946,rankin1961}; covering an interval by translates of such
sets already gives superpolynomial growth in the number of colors
\cite[pp.~2--3]{foxHunter2026}. For three-term progressions, the density
bound of Kelley and Meka gives a quasipolynomial coloring upper bound
\cite[Theorems~1.1--1.2]{kelleyMeka2023}. For each fixed length at least
five, Leng, Sah, and Sawhney obtained further density bounds, improving
the earlier general coloring upper bounds
\cite[Theorem~1.1]{lengSahSawhney2024}.

For completeness, Theorem~\ref{transfer:large-r} gives a direct elementary
lower bound with an explicit constant uniform over all lengths $k\ge3$.
Combining it with the coloring upper bound in the
companion paper~\cite[Equation~(1.3)]{openaiArithmeticProgressions2026}
gives, for each integer $k\ge3$, a constant $A_k\ge1$ such that
for every integer $r\ge256$,
\[
 \exp\!\left(\frac{(\log r)^2}{64\log2}\right)
 <W_r(k)\le
 \left\lceil\exp\bigl(A_k(2+\log r)^{A_k}\bigr)\right\rceil.
\]
Thus, for every fixed progression length $k\ge3$, as integer $r$ tends
to infinity, growth in the number of colors is superpolynomial but at
most quasipolynomial: the lower
bound gives $W_r(k)/r^D\to\infty$ for every fixed real $D>0$, while
the upper bound has logarithm at most a fixed power of $\log r$.
The upper constant and exponent may depend on $k$; no uniform bound
on that dependence or matching logarithmic exponent is asserted.
Appendix~\ref{sec:upper} gives an elementary recursive upper bound,
including a proof of finiteness.
This is a coarse version of the classical argument; quantitative upper
bounds include Shelah's primitive-recursive bound~\cite{shelah1988}
and the bounds following from Gowers's proof of Szemer\'edi's
theorem~\cite{gowers2001}.
The boundary cases are $W_1(k)=k$, $W_r(1)=1$, and $W_r(2)=r+1$.

\paragraph{The construction and its quantitative ingredient.}
We first construct a two-coloring of a cyclic group $\Z/q^D\Z$ with
$q^D\ge k^{ck}$, where $q$ is a power of a prime exceeding $k$.
Let $\lambda$ be an integer dilation, to be chosen below.
The homomorphisms
\[
 x_i(n)=n/q^i\pmod1,\qquad y_i(n)=\lambda x_i(n),
 \qquad 1\le i\le D,
\]
supply two systems of circle coordinates.
Partition each coordinate circle into small intervals, and label a point
by the tuple of intervals containing its coordinates.
An outer color is assigned to each such label.
The intervals for $y$ become finer near the cut used to select real
representatives.
We arrange that every $k$-term cyclic progression with nonzero step
either contains a positive proportion of each outer color or has
globally affine real representatives in the $y$ coordinates.

A heavily repeated label determines a short rational period $h$.
The quantitative obstacle is to control such returns when the dimension
grows with $k$: constants depending arbitrarily on the dimension would
not suffice.
The quantitative choice is to let $D$ grow and to make $\lambda$
contain the full primary part, at every prime at most $D^2$, of
every possible period $h\le2\lceil k^{1/2}\rceil$.
After dilation, each scalar rational period is either one or exceeds
$D^2$.
This separation permits a local test on entire period blocks.
A key estimate counts the patterns through one specified label by
$(CDh)^{6D}$, with an absolute $C$ independent of $q$, $\lambda$,
and the width of the anchor interval.
It follows by anchoring at one visit and counting affine comparisons
in two scalar parameters at a time.
This form of counting remains useful even when some intervals are
arbitrarily small.

We attach an independent Bernoulli$(p)$ flip to each key consisting
of a $y$-box and $\lfloor\|q\widetilde y\|_2^2\rfloor$, where
$p=k^{-1/20}$ and $\widetilde y\in[-1/2,1/2)^D$ is the centered
representative. Here $\|\cdot\|_2$ is the Euclidean norm.
Every occurrence of a key uses the same flip, and each key occurs
at most four times along a progression.
Making a progression rich in both outer colors monochromatic
therefore requires many rare flips.
In the possibly overlapping affine case, monochromaticity prescribes
values for many distinct flip variables. The global label-word count,
together with the three coefficients of the quadratic squared norm,
gives only a small family of words of outer colors and actual keys;
identical words prescribe the same event.
With $D=\lceil k^{1/10}\rceil$, their logarithmic count is
$O(k^{9/10}\log k)=o(pk)$.
The two resulting probability estimates eliminate every progression.
A deterministic digit product then gives all $r\ge2$.

The squared-norm geometry continues the tradition of
Behrend's construction~\cite{behrend1946}; Green and Wolf
\cite{greenWolf2010} give a torus-annulus proof of Elkin's refinement.
For the asymmetric problem of excluding a red three-term progression
and a blue long progression, Green~\cite{green2022} uses translates
of a randomly perturbed ellipsoidal annulus, while
Hunter~\cite{hunter2022} chooses annulus radii independently at many centers.
Here the randomness is attached to coordinate boxes and integer
squared-norm bands. Least-common-multiple dilations in cyclic groups
appear in Fox and Hunter~\cite{foxHunter2026}; digit carry control and
finite-field norm colorings for a related symmetric-progression problem
appear in Shi and Dong~\cite{shiDong2026}.
We prove all geometric, counting, and probabilistic ingredients needed
here, including the finite local lemma.

\paragraph{Organization.}
Section~\ref{sec:setup} fixes the cyclic group and the two meshes.
Section~\ref{sec:counting} proves the global and anchored pattern counts,
and Section~\ref{sec:outer} uses them to construct the outer coloring.
Section~\ref{sec:dichotomy} turns its local tests into the balanced-or-affine
alternative. Section~\ref{perturb:section} introduces the norm-indexed
perturbation and proves the cyclic two-color theorem.
Section~\ref{sec:transfers} transfers that theorem to integer intervals
and all color counts, and proves the additional growth statements.
Appendix~\ref{sec:upper} supplies the elementary finite upper bound.
